\subsection{Descompresión Selectiva y Estrategia de Acceso}

Una vez adoptada la representación por bloques, el principio operativo de la arquitectura es que solo deben descomprimirse los bloques necesarios
para ejecutar una operación. La descompresión selectiva constituye, por tanto, la segunda decisión de diseño y el mecanismo que vincula el ahorro
de memoria con el patrón real de acceso del simulador.

En un simulador tipo Schrödinger, muchas compuertas actualizan amplitudes cuya relación puede resolverse localmente dentro de una región acotada del
vector, mientras que otras inducen interacciones entre posiciones que pertenecen a bloques distintos. La arquitectura debe poder manejar ambos
casos sin perder consistencia. Cuando todas las amplitudes requeridas por la operación se encuentran en un mismo bloque, el acceso es enteramente
local: el bloque se materializa en memoria densa, se ejecuta la actualización y luego se decide si permanece temporalmente activo o si vuelve a
almacenamiento comprimido. Cuando la operación involucra amplitudes distribuidas en bloques distintos, la arquitectura debe adquirir
simultáneamente los bloques necesarios, aplicar la actualización y coordinar su posterior permanencia o escritura diferida.

Este enfoque permite que la compresión actúe de forma compatible con la lógica del simulador. El costo de acceso deja de depender del tamaño total
del vector y pasa a depender del número de bloques activados por cada compuerta, de la localidad del patrón de acceso y de la frecuencia con la que
esos mismos bloques vuelven a ser requeridos. En consecuencia, la efectividad de la arquitectura queda directamente vinculada a la estructura del
estado y al tipo de circuito ejecutado, lo cual resulta coherente con el análisis previo de compresibilidad.
